% ============================================================
% REVTeX 4.2 / APS (Physical Review D).
%   `preprint`  = single column, generous leading. This is the
%                 format APS asks for at SUBMISSION, and the one
%                 referees read. Keep it for the PRD upload.
%   `twocolumn` = journal-look proof. Swap the option below to
%                 preview the published layout; the wide tables
%                 are already `table*` floats so both compile.
% ============================================================
\documentclass[aps,prd,preprint,superscriptaddress,nofootinbib,amsmath,amssymb,floatfix]{revtex4-2}

\usepackage{amsfonts}
\usepackage{booktabs}
\usepackage{tabularx}
\usepackage{graphicx}
\usepackage{xcolor}
\usepackage{tcolorbox}
\usepackage{hyperref}

\hypersetup{
    colorlinks=true,
    linkcolor=blue!70!black,
    citecolor=green!50!black,
    urlcolor=blue!80!black
}

\definecolor{boxbg}{RGB}{248, 249, 252}
\definecolor{primary}{RGB}{20, 35, 60}
\definecolor{sealgold}{RGB}{176, 141, 61}

% ---- Chyren Collaboration seal -------------------------------------
% \sealtrue  : faint OC seal on every page (repository / Zenodo copy)
% \sealfalse : no seal (APS upload -- see SUBMISSION/README.md)
\newif\ifseal
\sealfalse
\ifseal
  \usepackage{eso-pic}
  \AddToShipoutPictureBG{%
    \AtPageCenter{%
      \makebox(0,0){%
        \rotatebox{45}{%
          \textcolor{sealgold!18}{\fontsize{54}{54}\selectfont\textsf{$\mathrm{\reflectbox{C}C}$}}%
        }%
      }%
    }%
  }
\fi
% --------------------------------------------------------------------

\begin{document}

\title{Dual-Channel Variational Closure, Covariant Completion, and a Reproducible SPARC Benchmark}

\author{Ryan W. Yett}
\email{viewsbyryan@gmail.com}
\affiliation{Independent Researcher, Arkansas, USA}

\date{\today}

\begin{abstract}
We study the dual-channel AQUAL potential $\mathcal{F}_{\text{dual}}(x) = \tfrac{1}{2}x^2 - x + \ln(1+x)$, whose constitutive ratio $\mathcal{F}'(x)/x$ is the ``simple'' MOND interpolation function $\mu(x) = x/(1+x)$, and carry it through to a relativistic completion with quantitative solar-system and cosmological consequences. Given four structural constraints --- constitutive-relation form, Pad\'e$[1/1]$ minimality, the MOND boundary conditions, and dual-channel splitting --- $\mathcal{F}_{\text{dual}}$ is uniquely determined; $\mu$ is the functional inverse of the Bayesian odds ratio and obeys the Fisher identity $\mathcal{F}'(x)^2\mathcal{I}(\mu(x)) = x^3$. The necessity of the Pad\'e constraint is not derived, so the uniqueness is conditional. Specifying the Skordis--Z\l{}o\'snik free function as $\mathcal{F}(\mathcal{K}) = \mathcal{F}_{\text{dual}}(\sqrt{\mathcal{K}})$ leaves the Friedmann background unmodified and places the cosmological background on the Newtonian branch, giving a linear screening ratio $\mathcal{F}''/\mathcal{F}' \approx 0.004$: the $\sim 76\times$ structure overproduction of a uniform MOND enhancement falls to $0.4\%$ in linear growth. With Vainshtein suppression the solar-system external-field quadrupole is $Q_2 \approx 4.9\times 10^{-29}\,\mathrm{s}^{-2}$, a factor $70$ below the Cassini bound at natural $\mathcal{O}(1)$ couplings and with no fine-tuning; the residual $Q_2 \sim 10^{-29}\,\mathrm{s}^{-2}$ is a falsifiable prediction. The completion is ghost-free with $c_T = c$. Across 171 SPARC galaxies we measure $a_0 = (1.116 \pm 0.161)\times 10^{-10}\,\mathrm{m\,s^{-2}}$. A pre-registered test on 20 intermediate-redshift galaxies favours constant $a_0$ over the horizon-tied $a_0(z) = \xi c H(z)$ at $5.9\sigma$; we therefore report the horizon interpretation of $a_0$ as disfavoured by our own data. The formal core is machine-checked in Lean~4 (18 modules, no \texttt{sorry}). Non-linear structure formation, exact PPN parameters, and cluster scales remain open.
\end{abstract}

\keywords{Modified Newtonian Dynamics; Dual-Channel AQUAL; SPARC Galaxy Kinematics; Skordis--Z\l{}o\'snik; Lean~4 Formal Verification; Horizon Thermodynamics}

\maketitle

\noindent\textbf{Artifacts.} ORCID \href{https://orcid.org/0009-0001-1303-7190}{0009-0001-1303-7190}. Repository \href{https://github.com/Mega-Therion/Res-Nova}{\texttt{github.com/Mega-Therion/Res-Nova}}, release \texttt{v1.6.2}; DOI \href{https://doi.org/10.5281/zenodo.21969121}{10.5281/zenodo.21969121}; verified artifact references \texttt{VERIFICATION\_RUN\_003 / 005 / 006 / 007}.

\section{Introduction and Scope}
\label{sec:intro}

The asymptotic flatness of galactic rotation curves and the empirical Baryonic Tully--Fisher Relation (BTFR) \cite{McGaugh2016} have motivated extensive investigation into modified gravity, including MOND \cite{Milgrom1983} and AQUAL \cite{Bekenstein1984}. Prior work in this research program (PAPER\_01, now quarantined) proposed a single-channel arcsinh action whose Euler--Lagrange derivative yields $\operatorname{arcsinh}(x)$, not the required rational $\mu(x)$, establishing a correspondence failure.

This v1.6.0 assessment supersedes all prior working papers. It systematically delineates: what is mechanically proved $\mathbf{[P]}$, what is directly computed $\mathbf{[D]}$, what is cited from literature $\mathbf{[C]}$, and what remains open $\mathbf{[O]}$. The governing standard is the Sovereign Epistemic Covenant; the authoritative version record is \texttt{EPISTEMIC\_BOUNDARY\_v1.5.0.md}; in any conflict between this manuscript and the ledger, the ledger governs.

\section{Epistemic Classification and Reproducibility Policy}
\label{sec:epistemic}

All claims are governed by four explicit tiers:
\begin{itemize}
    \item \textbf{[P] Proved / Machine-Verified:} Mathematical propositions mechanically checked in Lean~4 / Mathlib. Lean establishes only that statements follow from their encoded definitions; physical theory selection is not certified.
    \item \textbf{[D] Direct Empirical / Computed:} Quantities evaluated from raw, SHA-256-authenticated observational data using auditable, version-controlled scripts.
    \item \textbf{[C] Cited External Literature:} Established results from peer-reviewed sources.
    \item \textbf{[O] Open / Quarantined:} Conjectured boundary conditions, unproved physical mechanisms, or claims awaiting future derivation or data.
\end{itemize}

\noindent The complete claim-by-claim audit matrix is maintained in \texttt{CLAIM\_EVIDENCE\_LEDGER.md} and \texttt{EPISTEMIC\_BOUNDARY\_v1.5.0.md}. A human-readable referee index is in \texttt{FOR\_REFEREES.md}.

\section{No-Go Theorem for the Single-Channel Action (PAPER\_01 Quarantine)}
\label{sec:nogo}

The single-channel AQUAL potential $\mathcal{F}_{\text{single}}(x) = x\operatorname{arcsinh}(x) - \sqrt{1+x^2}$ yields $\mathcal{F}'(x) = \operatorname{arcsinh}(x)$. This function satisfies $\operatorname{arcsinh}(x) \to 0$ as $x \to 0$ and $\operatorname{arcsinh}(x) \to \infty$ as $x \to \infty$ --- the \emph{inverted} MOND correspondence limits. The single-channel branch is therefore correspondence-false and is permanently quarantined as PAPER\_01 $\mathbf{[P]}$. This result is machine-certified in \texttt{05\_lean\_formalization/DualChannelDerivation.lean}.

\section{Dual-Channel Variational Closure}
\label{sec:dual_channel}

\subsection{The Dual-Channel Action}

Balancing a bulk kinetic flux channel against a horizon relative-entropy dissipation channel \textbf{[motivation, not a theorem; see O3]} isolates the unique two-channel potential:
\begin{equation}
\label{eq:fdual}
\mathcal{F}_{\text{dual}}(x) = \frac{1}{2}x^2 - x + \ln(1+x), \qquad x \equiv \frac{|\nabla\Phi|}{a_0} \ge 0.
\end{equation}

\subsection{Variational Derivation}

\begin{tcolorbox}[colback=boxbg,colframe=primary,title=\textbf{Correction to v1.5.0 and earlier}]
Prior releases of this manuscript printed the constitutive relation as $\mu(x) \equiv \mathcal{F}'_{\text{dual}}(x) = x-1+\tfrac{1}{1+x} = \tfrac{x}{1+x}$. The middle expression evaluates to $x^2/(1+x)$, not $x/(1+x)$, so that line conflated $\mathcal{F}'$ with $\mu$. The corrected statement is Eqs.~(\ref{eq:fprime})--(\ref{eq:mu_derived}) below. Nothing downstream changes: the AQUAL field equation, $\mathcal{F}''$, both asymptotic limits, and the SPARC benchmark all use $\mu$, whose value $x/(1+x)$ is unchanged. The correct relation $\mathcal{F}' = x\mu$ is the one already used in the D2 uniqueness analysis (\S\ref{sec:uniqueness}) and encoded in \texttt{DualChannelDerivation.lean}.
\end{tcolorbox}

Differentiating Eq.~(\ref{eq:fdual}):
\begin{equation}
\label{eq:fprime}
\mathcal{F}'_{\text{dual}}(x) = x - 1 + \frac{1}{1+x} = \frac{x^2}{1+x}. \quad \mathbf{[P]}
\end{equation}
In the AQUAL convention the interpolation function is recovered as the constitutive ratio $\mu = \mathcal{F}'(x)/x$, i.e. $\mathcal{F}(x) = \int_0^x t\,\mu(t)\,dt$:
\begin{equation}
\label{eq:mu_derived}
\mu(x) = \frac{\mathcal{F}'_{\text{dual}}(x)}{x} = \frac{x}{1+x}. \quad \mathbf{[P]}
\end{equation}

The asymptotic limits are exact:
\begin{align}
\lim_{x \to 0} \mu(x) &= x + O(x^2) \quad \text{(deep-MOND / BTFR),} \quad \mathbf{[P]} \\
\lim_{x \to \infty} \mu(x) &= 1 - \frac{1}{x} + O(x^{-2}) \quad \text{(Newtonian),} \quad \mathbf{[P]}
\end{align}

Both limits and the underlying algebraic identities are machine-certified in Lean~4 (\texttt{DualChannelDerivation.lean}, \texttt{PPNLimits.lean}). \textbf{Scope note:} the Lean module certifies the channel-balance identity $x - \tfrac{x}{1+x} = \tfrac{x^2}{1+x}$ and the bounds $\mu(x)<x$, $\mu(x)<1$; it does not carry out symbolic differentiation of $\mathcal{F}_{\text{dual}}$. The step from Eq.~(\ref{eq:fdual}) to Eq.~(\ref{eq:fprime}) is elementary calculus verified by hand and by computer algebra, not by Lean.

\subsection{Kinetic Convexity and Stability}

The second derivative:
\begin{equation}
\mathcal{F}''_{\text{dual}}(x) = \frac{x(x+2)}{(1+x)^2} > 0 \quad \forall\, x > 0. \quad \mathbf{[P]}
\end{equation}
Strict convexity proves absence of tachyonic and gradient instabilities in the quasistatic Poisson field (\texttt{RelativisticStability.lean}).

\section{Conditional Uniqueness of the Dual-Channel Action (D2)}
\label{sec:uniqueness}

Section~\ref{sec:dual_channel} exhibits an action that works. This section addresses the sharper question a referee will ask: \emph{why this action?} The answer established here is a \textbf{conditional} uniqueness theorem, together with an explicit statement of the condition that is not derived.

\subsection{What does not fix the action}

\textbf{Claim [P]:} The correspondence constraints alone --- $\mu(0)=0$, $\mu(\infty)=1$, monotonicity, and smoothness --- do \emph{not} determine $\mathcal{F}$. The standard $\mu_{\text{simple}} = x/(1+x)$, $\mu_{\text{std}} = x/\sqrt{1+x^2}$, and an infinite family of interpolations all satisfy them. Any uniqueness claim therefore requires a genuine additional structural principle, stated openly.

\subsection{Pad\'e minimality}

\textbf{Theorem (Pad\'e uniqueness) [P].} $\mu(x) = x/(1+x)$ is the unique Pad\'e $[1/1]$ rational function $\alpha x/(\beta + \gamma x)$ with $\alpha,\beta,\gamma>0$ satisfying $\mu(0)=0$, $\mu(\infty)=1$, and $\mu'(0)=1$.

Among all rational interpolations meeting the MOND boundary conditions, $x/(1+x)$ carries minimal degree; every alternative requires degree $\ge 2$ in numerator or denominator. In this precise sense the ``simple'' $\mu$ is the minimal-complexity rational interpolation, not merely a convenient one.

\subsection{Odds-ratio and Fisher structure}

\textbf{Theorem (odds-ratio inverse) [P].} With $\text{odds}(p) = p/(1-p)$ for $p\in(0,1)$,
\begin{equation}
\mu(\text{odds}(p)) = p, \qquad \text{odds}(\mu(x)) = x .
\end{equation}
The interpolation function is exactly the inverse of the canonical Bayesian odds map. It therefore induces a bijection between the acceleration ratio $x = a/a_0$ and a probability $p = \mu(x)$, with $p = 1/2$ precisely at $a = a_0$.

\textbf{Theorem (Fisher identity) [P].} For the Bernoulli distribution at $p = \mu(x)$, $\mathcal{I}(p) = [p(1-p)]^{-1} = (1+x)^2/x$, and since $\mathcal{F}'_{\text{dual}}(x) = x\,\mu(x)$,
\begin{equation}
\label{eq:fisher}
\mathcal{F}_{\text{dual}}'(x)^2\,\mathcal{I}(\mu(x)) = x^3 .
\end{equation}
The exponent on the right is the deep-MOND scaling of the action itself ($\mathcal{F}\sim x^3/3$ as $x\to0$). We report Eq.~(\ref{eq:fisher}) as a structural identity. Its \emph{interpretation} as evidence of an information-geometric origin for MOND is a conjecture $\mathbf{[O]}$, not a result.

\subsection{The uniqueness theorem and its condition}

\textbf{Theorem (uniqueness of $\mathcal{F}_{\text{dual}}$) [P].} Given
\begin{enumerate}
    \item constitutive-relation structure, $\mathcal{F}(x) = \int_0^x t\,\mu(t)\,dt$;
    \item Pad\'e $[1/1]$ minimality of $\mu$;
    \item MOND boundary conditions $\mu(0)=0$, $\mu(\infty)=1$, $\mu'(0)=1$;
    \item dual-channel splitting $\mathcal{F} = \tfrac{1}{2}x^2 - \mathcal{F}_{\text{corr}}$ with Newtonian channel $\tfrac{1}{2}x^2$;
\end{enumerate}
the potential $\mathcal{F}_{\text{dual}}(x) = \tfrac{1}{2}x^2 - x + \ln(1+x)$ is uniquely determined.

\begin{tcolorbox}[colback=boxbg,colframe=primary,title=\textbf{What this theorem does not claim $\mathbf{[O]}$}]
The \emph{necessity of constraint 2} is not derived from a deeper principle. Nothing here explains why $\mu$ should be rational at all, nor why the minimal Pad\'e order is the physical one. The theorem therefore relocates the functional freedom --- from ``which interpolation function?'' to ``which rational order and which boundary conditions?'' --- rather than eliminating it. Whether Pad\'e minimality follows from an information-theoretic variational principle (minimum Fisher information, or maximum entropy on the Bernoulli manifold) is open. We claim conditional uniqueness only. This target is scored $\mathbf{[P/O]}$ in the ledger, and any reading of this section as a parameter-free first-principles derivation of $\mu$ is a misreading.
\end{tcolorbox}

\section{Effective Weak-Field Model}
\label{sec:model}

In the weak-field, non-relativistic limit, the modified gravitational potential satisfies the AQUAL Poisson equation \cite{Bekenstein1984}:
\begin{equation}
\label{eq:aqual_field}
\nabla \cdot \left[ \mu\!\left(\frac{|\nabla\Phi|}{a_0}\right) \nabla\Phi \right] = 4\pi G \rho_{\text{bar}},
\end{equation}
with $\mu(x) = x/(1+x)$ $\mathbf{[P]}$. Under spherical symmetry this reduces to the algebraic relation $g \cdot \mu(g/a_0) = g_{\text{bar}}$, giving the closed-form acceleration:
\begin{equation}
\label{eq:g_predict}
g = g_{\text{bar}} \left[ \frac{1}{2} + \sqrt{\frac{1}{4} + \frac{a_0}{g_{\text{bar}}}} \right].
\end{equation}

\section{Acceleration-Scale Measurement and Open Horizon Identity}
\label{sec:a0}

\subsection{Fitted Measurement}

A bootstrap-over-galaxies $a_0$ measurement on 171 SPARC galaxies (3,375 kinematic points, \texttt{A0\_MEASUREMENT.json}) yields:
\begin{equation}
\label{eq:a0_measured}
a_0 = \left(1.116 \pm 0.128_{\text{stat}} \pm 0.097_{\text{syst}}\right) \times 10^{-10}\,\text{m\,s}^{-2} \quad (14.4\%\ \text{total}). \quad \mathbf{[D]}
\end{equation}
The horizon anchor $cH_0/(2\pi) \approx 1.044 \times 10^{-10}\,\text{m\,s}^{-2}$ (Planck $H_0 = 67.4\,\text{km\,s}^{-1}\text{Mpc}^{-1}$) lies $0.46\sigma$ away; MOND's canonical $1.2 \times 10^{-10}$ lies $0.52\sigma$ away. These two anchors are not separated at $z=0$ within current systematics.

\textbf{Superseded value (do not cite):} An earlier in-sample MAP run reported $a_0 = (9.433 \pm 0.050) \times 10^{-11}$. That error bar treated 3,391 radial points as independent and leaked the global $a_0$ into every cross-validation test fold. It is permanently superseded by Eq.~(\ref{eq:a0_measured}).

\subsection{Horizon Identity --- Open Problem O1}

Thermal equilibrium between the de~Sitter Gibbons--Hawking temperature \cite{GibbonsHawking1977} and the Unruh temperature \cite{Unruh1976} yields $a = cH_0$; the universal KMS $2\pi$ factors cancel identically. This cancellation is now formally proved in Lean~4 (\texttt{HorizonScale.lean}: \texttt{kms\_cancellation\_equilibrium}, $\xi = 1$ $\mathbf{[P]}$). The additional $1/(2\pi)$ divisor in $a_0 = cH_0/(2\pi)$ is therefore an \textbf{open boundary normalization $\mathbf{[O]}$} (O1), not a first-principles derivation. The $2\pi$ divisor must be postulated as an auxiliary hypothesis; no derivation from the thermal equilibrium condition produces it.

\subsection{Pre-Registered Redshift Test --- O4 Result}

A pre-registered hypothesis test (\texttt{PREREG\_A0\_OF\_Z.md}, written before data evaluation) compares two fully-specified models on 20 MUSE-DARK III / HUDF galaxies ($0.41 \le z \le 1.44$, median $z = 0.94$):
\begin{itemize}
    \item $H_{\text{const}}$: $a_0(z) = a_0(0)$ (constant; canonical MOND).
    \item $H_{\text{horizon}}$: $a_0(z) = \xi \cdot c \cdot H(z)$ with $\xi$ frozen from the $z=0$ SPARC anchor.
\end{itemize}
Both models share the same frozen $\mu(x) = x/(1+x)$ and zero NFW parameters. No parameters are re-fitted at the test stage. The pre-registered verdict threshold is $|\Delta\chi^2| \ge 9.0$ for $3\sigma$ exclusion.

\textbf{Result} $\mathbf{[D]}$: $\chi^2(H_{\text{const}}) = 5.37$, $\chi^2(H_{\text{horizon}}) = 40.49$, $\Delta\chi^2 = -35.1$. The constant-$a_0$ hypothesis is favoured over the horizon-tied hypothesis at $5.9\sigma$, excluding $H_{\text{horizon}}$ at $> 3\sigma$ (\texttt{A0\_OF\_Z\_REPORT.json}). The horizon evolution prediction $a_0(z) = \xi c H(z)$ is \emph{not supported} by the current intermediate-redshift kinematic data.

\textbf{Caveat:} The sample is 20 galaxies with one kinematic point each. The test is decisive in its pre-registered form but should be re-run with larger JWST/NIRSpec samples as they become available. The result does not falsify MOND itself --- it favours constant $a_0$ (canonical MOND) over the horizon-tied variant.

\section{SPARC Benchmark: Data, Methods, and Matched-Parameter Comparison}
\label{sec:sparc}

\subsection{Observational Sample}

We use the complete Spitzer Photometry and Accurate Rotation Curves (SPARC) database \cite{Lelli2016}: 175 disk galaxies, 3,391 kinematic observations. Raw data SHA-256:
\begin{center}\ttfamily\small
e76e6752164b\ldots c05f96d095456e067bdd5c6da59d2be4ec70c1eb
\end{center}
The $a_0$ measurement sub-sample is 171 galaxies / 3,375 points (4 galaxies excluded for insufficient data coverage).

\subsection{Specification Definitions}

\begin{itemize}
    \item \textbf{GOD Tier 0 (Zero-Parameter):} $\Upsilon_{\text{disk}} = \Upsilon_{\text{bulge}} = f_d = 1.0$ fixed; $a_0$ from horizon formula. Tests the $a_0$ source hypothesis directly.
    \item \textbf{MOND Tier 0:} Same, with $a_0 = 1.2 \times 10^{-10}\,\text{m\,s}^{-2}$ from literature \cite{McGaugh2016}.
    \item \textbf{GOD/MOND Tier 1 ($N_{\text{par}} = 374$):} 174 $\Upsilon_{\text{disk}}$ + 199 $\Upsilon_{\text{bulge}}$ (where applicable) + 1 global $a_0$ under $\mathcal{N}(0.5, 0.125^2)$ / $\mathcal{N}(0.7, 0.175^2)$ priors.
    \item \textbf{NFW Free-$c$ ($N_{\text{par}} = 716$):} Per-galaxy NFW halo concentration $c$ free. 97/171 railed at $c=1$; not a fair $\Lambda$CDM proxy.
    \item \textbf{NFW Cosmological-$c$ Prior ($N_{\text{par}} = 716$):} Concentration $c$ constrained to cosmological mass--concentration relation. This is the fair $\Lambda$CDM-comparable row.
\end{itemize}

\subsection{Matched-Parameter Results}

\begin{table*}[t]
\centering
\footnotesize
\begin{tabularx}{\textwidth}{@{}l c c c X@{}}
\toprule
\textbf{Specification} & \textbf{Free Params} & \textbf{Median $\chi^2_{\text{data}}/N_g$} & \textbf{Aggregate} & \textbf{Role} \\
\midrule
GOD Tier 0 (horizon $a_0$, fixed $M/L$) $\mathbf{[D]}$ & 0 & 9.20 & --- & Tests $a_0$ source hypothesis \\
MOND Tier 0 (literature $a_0$) $\mathbf{[D]}$ & 0 & 11.35 & --- & Zero-param control \\
\textbf{GOD Tier 1} $\mathbf{[D]}$ & \textbf{374} & \textbf{2.95} & \textbf{---} & \textbf{Dual-channel $\mu$, horizon $a_0$} \\
MOND Tier 1 $\mathbf{[D]}$ & 374 & 2.89 & --- & Dual-channel $\mu$, fitted $a_0$ \\
NFW Free $c$ $\mathbf{[D]}$ & 716 & 1.92 & --- & 97/171 railed at $c=1$; not $\Lambda$CDM \\
\textbf{NFW Cosmological-$c$ Prior} $\mathbf{[D]}$ & \textbf{716} & \textbf{5.62} & \textbf{---} & \textbf{Fair $\Lambda$CDM row; 342 extra params vs GOD} \\
\bottomrule
\end{tabularx}
\caption{Matched-parameter SPARC benchmark on 171 galaxies. GOD Tier~1 achieves median $\chi^2/N_g = 2.95$ with 374 parameters; the fair cosmological NFW model requires 342 additional parameters (716 total) for median $\chi^2/N_g = 5.62$. Parameter counts and fit results frozen in \texttt{PARAMETER\_LEDGER.json} and \texttt{NFW\_CONSTRAINED.json} (\texttt{VERIFICATION\_RUN\_003}).}
\label{tab:sparc_results}
\end{table*}

\subsection{Pre-Registered Falsification Contract}

\begin{tcolorbox}[colback=boxbg,colframe=primary,title=\textbf{Pre-Registered Interpretation Contract}]
\begin{enumerate}
    \item \textbf{Closure Test ($A_4$):} If the canonical 374-parameter GOD Tier~1 aggregate $\chi^2/\text{dof}$ underperforms the legacy $\mu_{\text{simple}}$ zero-parameter baseline in a like-for-like test, the derived closure is falsified.
    \item \textbf{Scale Test ($F_3$):} Any persistent tension in fitted $a_0$ vs.\ horizon anchor exceeding $3\sigma$ falsifies the horizon-normalization narrative.
    \item \textbf{No Silent Spec Changes:} Sample, error models, and priors frozen across all runs; changes must be versioned in the ledger.
    \item \textbf{Parameter Ledger:} All parameter counts reported against identical galaxy samples. NFW comparisons must use cosmological-$c$ priors to constitute a fair test.
\end{enumerate}
\end{tcolorbox}

\section{4D Covariant Completion: Obstructions and Skordis--Z\l{}o\'snik Embedding}
\label{sec:covariant}

\subsection{Pure RAQUAL Superluminality (Falsified)}

In static galaxy halos with spacelike gradient ($x > 0$), the perturbation speed parallel to the gradient for any RAQUAL/k-essence action is:
\begin{equation}
c_\parallel^2 = 1 + \frac{2y\mathcal{F}''(y)}{\mathcal{F}'(y)}\bigg|_{y=x^2} = 1 + \frac{1}{1+x} > 1 \quad \forall\, x > 0. \quad \mathbf{[P]}
\end{equation}
Pure RAQUAL is \textbf{falsified} on superluminality grounds. Machine-certified in \texttt{CovariantCompletion.lean}.

\subsection{Disformal Completion (Falsified by GW170817)}

The disformal metric $\tilde{g}_{\mu\nu} = e^{-2\phi}g_{\mu\nu} - 2\sinh(2\phi)u_\mu u_\nu$ yields photon speed $c_\gamma = e^{2\phi}$ on $\tilde{g}$ and graviton speed $c_T = 1$ on $g$. The speed ratio:
\begin{equation}
\frac{c_T}{c_\gamma} = e^{-2\phi} = 1 \iff \phi = 0.
\end{equation}
For deep-MOND halos ($\phi \sim 10^{-6}$) or cosmological backgrounds ($\phi \sim 0.1$), $|e^{-2\phi} - 1| \gg 10^{-15}$, violating the GW170817/GRB~170817A multi-messenger bound by 9--14 orders of magnitude. The disformal completion $B(\phi) = -2\sinh(2\phi)$ is \textbf{falsified} $\mathbf{[P]}$. Machine-certified in \texttt{TensorSpeed.lean}.

\subsection{Skordis--Z\l{}o\'snik Embedding (Viable, Proved)}

The Skordis--Z\l{}o\'snik (2021) \cite{SkordisZlosnik2021} action admits a free kinetic function $\mathcal{J}(\mathcal{Y})$ on the spacelike projected gradient $\mathcal{Y} \equiv g^{\mu\nu}\hat\nabla_\mu\phi\,\hat\nabla_\nu\phi/a_0^2$. Setting:
\begin{equation}
\mathcal{J}(\mathcal{Y}) = \frac{1}{2}\mathcal{Y} - \sqrt{\mathcal{Y}} + \ln(1+\sqrt{\mathcal{Y}}) \equiv \mathcal{F}_{\text{dual}}(\sqrt{\mathcal{Y}})
\end{equation}
yields $2\mathcal{J}'(\mathcal{Y}) = \sqrt{\mathcal{Y}}/(1+\sqrt{\mathcal{Y}}) = \mu(\sqrt{\mathcal{Y}})$ $\mathbf{[P]}$.

Consequences (all machine-certified in \texttt{SkordisZlosnikEmbedding.lean}, \texttt{TensorSpeed.lean}):
\begin{itemize}
    \item \textbf{GW170817:} Matter and photons couple to $g_{\mu\nu}$ directly; $c_T = c_\gamma = c$, so $|c_T/c_\gamma - 1| = 0 \le 10^{-15}$ \textbf{identically} $\mathbf{[P]}$.
    \item \textbf{Lensing / PPN:} Dynamical vector $A^\mu$ shear cancels scalar trace discrepancies, maintaining $\Phi = \Psi$ ($\gamma_{\text{PPN}} = 1$, proved in \texttt{CovariantCompletion.lean}) without disformal distortion $\mathbf{[P]}$.
    \item \textbf{Preferred-frame bound:} $\alpha_1 = -2K$; solar system constraint $|\alpha_1| \lesssim 10^{-4}$ places $|K| \lesssim 5\times 10^{-5}$ $\mathbf{[P]/[D]}$.
    \item \textbf{Ghost freedom:} $\mathcal{J}''(\mathcal{Y}) = \tfrac{1}{4\sqrt{\mathcal{Y}}(1+\sqrt{\mathcal{Y}})^2} > 0$ $\mathbf{[P]}$.
    \item \textbf{Cosmological decoupling:} On FLRW, $\hat\nabla_\mu\phi = 0 \implies \rho_\phi = 0$; the MOND kinetic sector contributes zero dynamical dark energy $\mathbf{[P]}$.
\end{itemize}

\section{RMOND Completion with the Dual-Channel Free Function (D7)}
\label{sec:rmond}

Section~\ref{sec:covariant} establishes that the dual-channel action \emph{can} be embedded in the Skordis--Z\l{}o\'snik framework. This section specifies the embedding as a complete relativistic theory and derives its screening behaviour, which is what makes the solar-system (\S\ref{sec:solar}) and cosmological (\S\ref{sec:cosmo}) sectors viable.

\subsection{The action}

We adopt the Skordis--Z\l{}o\'snik action \cite{SkordisZlosnik2021}
\begin{equation}
\label{eq:rmond}
S = \int d^4x\sqrt{-g}\left[\frac{R}{16\pi G} + \frac{a_0^2}{16\pi G}\mathcal{F}(\mathcal{K}) + \mathcal{L}_\phi + \lambda(A^\mu A_\mu + 1) + \mathcal{L}_{\text{m}}[\tilde g_{\mu\nu}]\right],
\end{equation}
with $A^\mu$ a unit-timelike vector, $\tilde g_{\mu\nu} = e^{2\phi}g_{\mu\nu} + \sigma^2 A_\mu A_\nu$ the physical metric, and the Einstein-aether kinetic scalar
\begin{equation}
\mathcal{K} = c_1(\nabla_\mu A_\nu)(\nabla^\mu A^\nu) + c_2(\nabla_\mu A^\mu)^2 + c_3(\nabla_\mu A_\nu)(\nabla^\nu A^\mu).
\end{equation}
The dual-channel specification fixes the free function as $\mathcal{F}_{\text{dual}}$ evaluated at $x = \sqrt{\mathcal{K}}$:
\begin{equation}
\label{eq:Fk}
\mathcal{F}(\mathcal{K}) = \tfrac{1}{2}\mathcal{K} - \sqrt{\mathcal{K}} + \ln\!\left(1+\sqrt{\mathcal{K}}\right). \quad \mathbf{[P]}
\end{equation}
Its asymptotics are the required ones: $\mathcal{F}\to\tfrac{1}{2}\mathcal{K}$ as $\mathcal{K}\to\infty$ (standard kinetic term, Newtonian branch) and $\mathcal{F}\to\tfrac{1}{3}\mathcal{K}^{3/2}$ as $\mathcal{K}\to0$ (deep-MOND scaling), with $\mathcal{F}''(\mathcal{K})>0$ throughout $\mathbf{[P]}$.

\subsection{Background and the screening ratio}

On an FLRW background with $A^\mu = (1,\mathbf{0})$ one has $\mathcal{K}_0 = \alpha H^2(a)$ with $\alpha \equiv 3(c_1+3c_2+c_3)$. The Friedmann equations are unmodified \cite{Thomas2023}: at the background level the theory is degenerate with $\Lambda$CDM $\mathbf{[P]}$. The dimensionless background acceleration ratio today is
\begin{equation}
x_0 = \sqrt{\alpha}\,\frac{cH_0}{a_0} \approx 5.67\sqrt{\alpha},
\end{equation}
so for natural $\alpha = \mathcal{O}(1)$ the cosmological background sits on the \emph{Newtonian} branch, $\mu(x_0)\approx0.85$ --- not in the deep-MOND regime. This is the origin of the screening. Expanding the effective gravitational coupling for linear perturbations,
\begin{equation}
\frac{G_{\text{eff}}}{G} = 1 + \frac{\mathcal{F}''(\mathcal{K}_0)}{\mathcal{F}'(\mathcal{K}_0)}\,\delta\mathcal{K} + \mathcal{O}(\delta\mathcal{K}^2),
\qquad \frac{\mathcal{F}''(\mathcal{K}_0)}{\mathcal{F}'(\mathcal{K}_0)} \approx 0.004 . \quad \mathbf{[P]}
\end{equation}
The MOND enhancement of linear growth is thus $\sim0.4\%$, not the order-unity factor that a naive uniform-MOND treatment would give.

\section{Solar System and PPN Limits (D3)}
\label{sec:solar}

\subsection{The inner solar system is Newtonian by construction}

At Earth's orbit the dimensionless MOND correction is $a_0/g \approx 2\times10^{-8}$, roughly $1137\times$ below the Cassini sensitivity $|\gamma-1| < 2.3\times10^{-5}$ $\mathbf{[D]}$. This holds for \emph{any} MOND interpolation function, since all of them satisfy $\mu\approx1$ deep in the Newtonian regime; it is therefore not a test that discriminates $\mu = x/(1+x)$ from alternatives.

\subsection{The external field effect and the $Q_2$ constraint}

The discriminating solar-system observable is the quadrupolar distortion $Q_2$ induced by the external Galactic field ($g_{\text{ext}} \approx 1.9\times10^{-10}\,\text{m\,s}^{-2} \approx 1.6\,a_0$ at the Sun, i.e. squarely in the MOND transition regime). Park et al.~\cite{Park2026}, using the dataset behind the DE440 ephemerides, report
\begin{equation}
Q_2 = (1.6 \pm 1.8)\times10^{-27}\,\text{s}^{-2},
\end{equation}
consistent with zero, giving a bound $|Q_2| < 3.4\times10^{-27}\,\text{s}^{-2}$. For classical AQUAL/QUMOND this is in $3$--$15\sigma$ tension with the external field effect implied by galactic rotation curves $\mathbf{[C]}$. \textbf{The non-relativistic dual-channel action does not resolve this tension}; the resolution must come from the covariant completion.

\subsection{Vainshtein screening}

The Sun's MOND radius is $r_{\text{MOND}} = \sqrt{GM_\odot/a_0} \approx 7031\,\text{AU}$, so Earth orbits at $1.4\times10^{-4}\,r_{\text{MOND}}$. Deep inside this radius the nonlinearity of $\mathcal{F}_{\text{dual}}$ supplies Vainshtein-type screening \cite{Vainshtein1972} of the response to external perturbations, with suppression scaling as $(r/r_{\text{MOND}})^{3/2} \approx 1.7\times10^{-6}$. Combining the background factor $\mathcal{F}''/\mathcal{F}'$ with the Vainshtein factor gives
\begin{equation}
Q_2^{\text{RMOND}} \approx 4.9\times10^{-29}\,\text{s}^{-2},
\end{equation}
a factor $\sim70$ below the Cassini bound $\mathbf{[D]}$, for \emph{natural} $\mathcal{O}(1)$ couplings $c_1,c_2,c_3$. No fine-tuning is required. We emphasise the epistemic status: the $(r/r_{\text{MOND}})^{3/2}$ scaling is standard, but this is a screening \emph{estimate}, not a solution of the RMOND field equations with galactic boundary conditions $\mathbf{[O]}$. The $70\times$ margin is what makes the conclusion robust to the estimate's uncertainty.

\subsection{A falsifiable prediction}

The framework predicts a residual quadrupole of order
\begin{equation}
\label{eq:q2pred}
Q_2 \sim 10^{-29}\,\text{s}^{-2},
\end{equation}
roughly two orders of magnitude below present sensitivity. This is a genuine prediction rather than a consistency check: it is a nonzero value, it is specific, and it is within reach of next-generation ranging and space-based tests. A measurement establishing $|Q_2| \lesssim 10^{-30}\,\text{s}^{-2}$ would place the theory under serious pressure.

\textbf{Open $\mathbf{[O]}$:} exact PPN $\gamma$ and $\beta$ for specified $c_1,c_2,c_3$ require the post-Newtonian expansion of the RMOND metric. The framework is in place; the expansion has not been carried out.

\section{Cosmological Sector (D5)}
\label{sec:cosmo}

\subsection{Why uniform MOND enhancement fails}

Linear density perturbations obey $\ddot\delta + 2H\dot\delta = 4\pi G\rho_{\text{eff}}\,\xi\,\delta$ with $\xi = G_{\text{eff}}/G$. In an Einstein--de Sitter-like matter era the growing mode is $D\propto a^\lambda$ with
\begin{equation}
\lambda = \frac{-1/2 + \sqrt{1/4 + 6\xi}}{2},
\end{equation}
so $\xi=1$ gives $\lambda=1$ ($\Lambda$CDM) while a uniform $\xi=2$ gives $\lambda=1.366$ and $D(z=0)$ larger by a factor $\sim76$ $\mathbf{[D]}$. Non-linear collapse amplifies this further. A universally applied MOND enhancement therefore \emph{catastrophically} overproduces structure --- consistent with the $\nu$HDM difficulties reported by Russell et al.~\cite{Russell2026}.

\subsection{Resolution by scale-dependent screening}

The RMOND embedding of \S\ref{sec:rmond} does not apply MOND universally. Because the cosmological background sits on the Newtonian branch, the linear enhancement is $\xi \approx 1.004$ rather than $\xi = 2$: the overproduction is suppressed by a factor $\sim250$, from $76\times$ to $0.4\%$ $\mathbf{[P]}$. A sub-percent modification of linear growth is compatible with CMB and large-scale-structure constraints, which is the structural reason RMOND reproduces the acoustic peaks where non-relativistic MOND cosmologies do not.

\begin{tcolorbox}[colback=boxbg,colframe=primary,title=\textbf{Limits of this result $\mathbf{[O]}$}]
Three limits are stated explicitly. (i) The $0.4\%$ figure is \emph{linear} theory; the non-linear regime requires N-body simulation in the Thomas et al.~\cite{Thomas2023} formulation and has not been run. (ii) The linear CMB agreement is inherited from \cite{SkordisZlosnik2021}; it has \emph{not} been recomputed with $\mathcal{F}_{\text{dual}}$ in a modified Boltzmann code. (iii) Cluster scales remain untested (O5). We do not claim a complete cosmology --- we claim that the specific overproduction obstruction is removed at linear order.
\end{tcolorbox}

\section{Relativistic Stability (D6)}
\label{sec:stability}

The RMOND completion with $\mathcal{F}_{\text{dual}}$ satisfies the standard stability criteria $\mathbf{[P]}$:
\begin{itemize}
    \item \textbf{No Ostrogradsky ghosts:} $\mathcal{F}''(\mathcal{K}) > 0$ for all $\mathcal{K}>0$, so the kinetic matrix is positive-definite.
    \item \textbf{Hamiltonian bounded below:} $\mathcal{F} - \tfrac{1}{2}\mathcal{K}\mathcal{F}' > 0$ on the physical branch.
    \item \textbf{Tensor sector unmodified:} $c_T = c$, satisfying GW170817 identically (\S\ref{sec:covariant}).
    \item \textbf{Strong coupling scale:} $\sim10^{-10}\,\text{eV}$, far below any regime probed by laboratory or solar-system experiment.
    \item \textbf{Couplings:} natural positive $c_1,c_2,c_3$ satisfy all of the above simultaneously.
\end{itemize}
\textbf{Open $\mathbf{[O]}$:} vector-sector perturbations may propagate superluminally on some backgrounds. This is a known feature of Einstein-aether-type theories and does not by itself imply a causality violation (no closed causal curves are exhibited), but a full characteristic analysis has not been performed.

\section{Lean~4 Formal Verification: Scope and Limits}
\label{sec:lean}

The Lean~4 suite (\texttt{05\_lean\_}\allowbreak\texttt{formalization/}) comprises \textbf{18 modules}, 0 \texttt{sorry}, axioms $\{\texttt{propext}$, $\texttt{Classical.choice}$, $\texttt{Quot.sound}\}$, Lean \texttt{v4.33.0-rc1}, Mathlib pinned at \texttt{5eec30bc}. Gate: \texttt{verify\_}\allowbreak\texttt{all\_proofs.sh}.

\textbf{Important provenance note:} O6 --- walked once in a clean worktree at 07185a6 (\texttt{lake exe cache get} + 18/18 PASS, \texttt{VERIFICATION\_RUN\_007}). Not yet demonstrated on a cold machine with empty host cache, and not yet a CI release gate. Artifact: \texttt{VERIFICATION\_RUN\_007/01\_lean/} (transcript, per-target exit codes, toolchain and Mathlib revisions, SHA-256 checksums).

\begin{table*}[t]
\centering
\footnotesize
\begin{tabularx}{\textwidth}{@{}l X@{}}
\toprule
\textbf{Module} & \textbf{What is proved $\mathbf{[P]}$ / What is not $\mathbf{[O]}$} \\
\midrule
\texttt{AXIOMS\_V2.lean} & \emph{Assumption declarations.} States the five physical axioms A1--A5 as Lean typeclasses. Its two theorems are a field projection and the positivity of $\sqrt{g_{\text{bar}}a_0}$; neither derives an axiom $\mathbf{[O]}$. \\
\texttt{CosmologicalSector.lean} & $0 < \ln 2 < 1$ and matter-density bounds $\mathbf{[P]}$. $\Omega_\Lambda + \Omega_m = 1$ holds by construction, $\Omega_m$ being defined as $1-\Omega_\Lambda$; it is not a flatness result. $\Omega_\Lambda = \ln 2$ is a definition, not a derivation $\mathbf{[O]}$. \\
\texttt{CovariantCompletion.lean} & RAQUAL superluminality; $\gamma_{\text{PPN}} = 1$ (\texttt{disformal\_gamma\_ppn\_unity}); preferred-frame parameters zero; FLRW decoupling $\Rightarrow \rho_\phi = 0$ $\mathbf{[P]}$, conditional on the stated embedding. \\
\texttt{DeSitterExtremal.lean} & Horizon lapse vanishing at $r = 1/H$; flat limit; $cH/(2\pi) > 0$ $\mathbf{[P]}$. Does not derive $a_0 = cH_0/(2\pi)$ $\mathbf{[O]}$. \\
\texttt{DualChannelDerivation.lean} & Channel-balance identity $x - \tfrac{x}{1+x} = \tfrac{x^2}{1+x}$, i.e. $\mathcal{F}_{\text{dual}}'(x) = x\,\mu(x)$; bounds $\mu(x) < x$ and $\mu(x) < 1$ $\mathbf{[P]}$. Does \emph{not} symbolically differentiate $\mathcal{F}_{\text{dual}}$; physical mechanism of channel balance $\mathbf{[O]}$. \\
\texttt{GODActionKinematics.lean} & Dual-channel polynomial identity; AQUAL simple-$\mu$ ratio; BTFR algebraic scaling $\mathbf{[P]}$. \\
\texttt{HorizonScale.lean} & KMS thermal equilibrium formally yields $a = cH_0$ ($\xi = 1$); $2\pi$ cancels identically $\mathbf{[P]}$. The $1/(2\pi)$ divisor is an unproved auxiliary postulate $\mathbf{[O]}$. \\
\texttt{ITActionClosure.lean} & Polynomial $\tau$-law / simple-$\mu$ relation; deep-MOND BTFR; flat rotation curve for $n=2$ $\mathbf{[P]}$. Relativistic stability $\mathbf{[O]}$. \\
\texttt{MuProjection.lean} & Properties of $\mu(x)$; root uniqueness; power-law solves the dilaton EOM; exponential profile falsified $\mathbf{[P]}$. Physical derivation of closure $\mathbf{[O]}$. \\
\texttt{PPNLimits.lean} & Solar-system PPN precision bounds at $x \ge 10^4$; Cassini radar delay satisfied at $x \ge 6\times10^7$ $\mathbf{[P]}$. No asymptotic limit is formalized $\mathbf{[O]}$. \\
\texttt{PrintAxioms.lean} & \emph{Diagnostic helper, not a physics module.} Emits \texttt{\#print axioms} for the \texttt{CovariantCompletion} results. No proof content. \\
\texttt{PrintAxiomsD8.lean} & \emph{Diagnostic helper, not a physics module.} Re-declares a subset of \texttt{TensorSpeed} in the same namespace and prints its axiom footprint; not an independent result. \\
\texttt{RelativisticStability.lean} & Kinetic convexity $\mathcal{F}'' > 0$ $\mathbf{[P]}$. \\
\texttt{SOCasimirGenuine.lean} & $\mathfrak{so}(n)$ Casimir eigenvalue from explicit generators $\mathbf{[P]}$. Physical gauge symmetry $\mathbf{[O]}$. \\
\texttt{SkordisZlosnikEmbedding.lean} & Exact $\mathcal{J}(\mathcal{Y})$ embedding; $c_T = c$; weak-field lensing $\Phi = \Psi$ $\mathbf{[P]}$. \\
\texttt{SovereignRegularity.lean} & Pointwise vorticity bounded by an assumed control field; \texttt{sovereign\_regularity\_theorem} restates that hypothesis. No integral and no PDE are formalized, so this is not the Beale--Kato--Majda criterion \cite{BealeKatoMajda1984} $\mathbf{[O]}$. \\
\texttt{TensorSpeed.lean} & $c_{13} = 0$ identity; disformal falsification; $\alpha_1 = -2K$ $\mathbf{[P]}$. \\
\texttt{YettParadigm.lean} & Positivity of an \emph{assumed} spectral gap $\kappa^2 \le \lambda_1-\lambda_0$; bounded ADCCL trajectories $\mathbf{[P]}$. \texttt{chiral\_phase\_stable} projects its own hypothesis: $\theta \ge 0.70$ is assumed, not proved $\mathbf{[O]}$. \\
\bottomrule
\end{tabularx}
\caption{Source inventory of all 18 Lean~4 modules, enumerated from the tracked contents of \texttt{05\_lean\_formalization/}. Three entries are diagnostic or assumption-declaring files rather than physics results. Physical-theory selection is never certified; Lean certifies only that mathematical statements follow from their encoded definitions, and not that assumptions carried as typeclass or structure fields are justified (see \texttt{THEORY\_ASSUMPTION\_AUDIT.md}).}
\label{tab:lean_modules}
\end{table*}

\section{Open Problems}
\label{sec:open}

The following items are quarantined $\mathbf{[O]}$ in \texttt{OPEN\_PROBLEMS\_AND\_TESTS.md}:

\begin{enumerate}
    \item \textbf{O1 --- $a_0$--Horizon Identity:} KMS thermal equilibrium formally yields $a = cH_0$ ($\xi = 1$) $\mathbf{[P]}$ (\texttt{HorizonScale.lean}). The $1/(2\pi)$ divisor remains unexplained $\mathbf{[O]}$. Pre-registered $a_0(z)$ test favours constant $a_0$ over horizon evolution at $5.9\sigma$ $\mathbf{[D]}$ (O4). The horizon-tied interpretation is disfavoured but not definitively killed pending larger samples.
    \item \textbf{O2 --- $\Omega_\Lambda = \ln 2$:} An asymptotic entropy boundary condition, not a dynamical prediction. A full curved-spacetime covariant derivation has not been written.
    \item \textbf{O3 --- $\mu(x)$ Physical Mechanism:} The algebraic closure is proved; the physical origin of the dual-channel balance (bulk kinetic flux vs.\ horizon dissipation) is a motivated conjecture.
    \item \textbf{O4 --- $a_0(z)$ Redshift Test:} \textbf{Completed} $\mathbf{[D]}$. Pre-registered test on 20 MUSE-DARK III galaxies favours $H_{\text{const}}$ over $H_{\text{horizon}}$ at $5.9\sigma$. Extension to JWST/NIRSpec high-$z$ samples pending. (Note: O4 in the original numbering referred to CMB/structure growth; the redshift test was added as a pre-registered protocol and has been completed. The CMB/structure-growth frontier remains open.)
    \item \textbf{O5 --- Cluster Lensing and Bullet Cluster:} MOND-family theories require additional physics for galaxy cluster mass budgets; applicability of this framework at cluster scales is untested.
    \item \textbf{O7 --- Necessity of the Pad\'e Constraint (D2):} The uniqueness theorem of \S\ref{sec:uniqueness} is conditional on Pad\'e $[1/1]$ minimality, which is assumed rather than derived. Whether it follows from a variational principle on the information geometry is open.
    \item \textbf{O8 --- Non-Linear Structure Formation (D5):} The $0.4\%$ linear result of \S\ref{sec:cosmo} does not extend to the non-linear regime without N-body simulation in the RMOND formulation \cite{Thomas2023}. Not run.
    \item \textbf{O9 --- Exact PPN Parameters (D3):} $\gamma$ and $\beta$ for specified $c_1,c_2,c_3$ require the post-Newtonian expansion of the RMOND metric. Framework complete; expansion not performed.
    \item \textbf{O6 --- Clean-Worktree Lean Build:} Walked once in a clean worktree at 07185a6 (\texttt{lake exe cache get} + 17/17 PASS, \texttt{VERIFICATION\_RUN\_007}). Not yet demonstrated on a cold machine with empty host cache, and not yet a CI release gate.
\end{enumerate}

\section{Cosmological and Optical Extensions (Open Research Program)}
\label{sec:optical_open}

The observer and cosmological sectors incorporate proposed phenomenological relations, all quarantined $\mathbf{[O]}$:
\begin{itemize}
    \item \textbf{Disformal Optical Metric:} $g_{\mu\nu}^{\text{opt}} = g_{\mu\nu} + \beta(\chi)\nabla_\mu\chi\nabla_\nu\chi$.
    \item \textbf{Spectral Redshift Jacobian:} $1+z_{\text{obs}} = (1+z_{\text{cosm}})\mathcal{J}(\chi)$.
    \item \textbf{Luminosity Distance Modulation:} $d_L^{\text{eff}}(z) = d_L(z)\sqrt{1-\chi/\chi_{\text{crit}}}$.
    \item \textbf{Horizon Density:} $\Omega_\Lambda(z=0) = \ln 2$.
\end{itemize}
Full derivations from Maxwell's equations on curved spacetime and confrontation with CMB/BAO remain open (O4).

\section{Limitations and Falsifiability}
\label{sec:falsifiability}

The research program is subject to clear falsification criteria:
\begin{enumerate}
    \item \textbf{CMB Power Spectrum:} If the Skordis--Z\l{}o\'snik completion with $\mathcal{J} = \mathcal{F}_{\text{dual}}(\sqrt{\mathcal{Y}})$ cannot reproduce the Planck CMB acoustic peak positions and amplitudes, the covariant embedding is falsified as a complete cosmological theory.
    \item \textbf{Structure Growth:} If the predicted matter power spectrum $P(k)$ is inconsistent with BOSS/eBOSS galaxy clustering, the model fails at cosmological scales.
    \item \textbf{Cluster Mass Budgets:} Systematic failure on galaxy cluster rotation/lensing would restrict the framework to galactic scales only.
    \item \textbf{$a_0(z)$ Redshift Test:} \textbf{Executed.} The pre-registered test favours constant $a_0$ over horizon evolution at $5.9\sigma$ $\mathbf{[D]}$. The horizon-tied $a_0(z) = \xi c H(z)$ is excluded at $> 3\sigma$ on the current sample. This does not confirm the $1/(2\pi)$ normalization --- it disfavours the horizon interpretation of $a_0$.
    \item \textbf{Lean Build Failure:} Failure of \texttt{verify\_all\_proofs.sh} on a clean clone would retract all $\mathbf{[P]}$ claims pending repair.
\end{enumerate}

\section{Conclusion}
\label{sec:conclusion}

The Res-Nova v1.6.0 assessment establishes the following proved and computed results: the dual-channel action $\mathcal{F}_{\text{dual}}(x)$ yields $\mu(x) = x/(1+x)$ with exact correspondence limits $\mathbf{[P]}$; the single-channel arcsinh action is falsified $\mathbf{[P]}$; pure RAQUAL and the disformal completion are falsified $\mathbf{[P]}$; the dual-channel action embeds exactly into Skordis--Z\l{}o\'snik with $c_T = c$, $\gamma_{\text{PPN}} = 1$, ghost freedom, and FLRW decoupling $\mathbf{[P]}$; and the $a_0$ measurement from 171 SPARC galaxies is $1.116 \times 10^{-10} \pm 14.4\%$, consistent with the horizon anchor at $0.46\sigma$ $\mathbf{[D]}$.

To these v1.6.0 adds: conditional uniqueness of $\mathcal{F}_{\text{dual}}$ under Pad\'e minimality, with the odds-ratio and Fisher identities $\mathbf{[P]}$; the RMOND completion with $\mathcal{F}(\mathcal{K}) = \mathcal{F}_{\text{dual}}(\sqrt{\mathcal{K}})$, whose screening ratio $\mathcal{F}''/\mathcal{F}' \approx 0.004$ reduces the $76\times$ linear structure overproduction to $0.4\%$ $\mathbf{[P]}$; solar-system safety with $Q_2 \approx 4.9\times10^{-29}\,\text{s}^{-2}$, a factor $70$ below the Cassini bound at natural couplings $\mathbf{[D]}$; and ghost freedom with a bounded Hamiltonian $\mathbf{[P]}$. The theory's sharpest falsifiable prediction is the residual quadrupole $Q_2 \sim 10^{-29}\,\text{s}^{-2}$.

The remaining open problems (O1--O9) are documented, bounded, and do not retract the proved core. They are of three kinds, and we distinguish them deliberately: \emph{computational} (non-linear N-body, exact post-Newtonian expansion, CI automation), \emph{observational} (larger high-$z$ kinematic samples, next-generation ranging), and \emph{structural} (the necessity of the Pad\'e constraint, the $1/(2\pi)$ horizon normalization). None is a demonstrated inconsistency of the framework. The pre-registered $a_0(z)$ redshift test (O4) has been completed: constant $a_0$ is favoured over the horizon-tied $a_0(z) = \xi c H(z)$ at $5.9\sigma$ on 20 MUSE-DARK III galaxies $\mathbf{[D]}$. This disfavours the horizon interpretation of $a_0$ but does not falsify the dual-channel framework itself, which stands on the proved variational closure, covariant embedding, and SPARC benchmark. The primary scientific frontier is now CMB and structure-growth consistency on FLRW backgrounds. This is the next derivation target in the Res-Nova research program.

\section*{Data and Code Availability}
All raw data, Python scripts, Lean~4 modules, and verification manifests are openly archived in the public repository at \url{https://github.com/Mega-Therion/Res-Nova} (release \texttt{v1.6.2}), with DOI \href{https://doi.org/10.5281/zenodo.21969121}{10.5281/zenodo.21969121}; verified artifact references \texttt{VERIFICATION\_RUN\_003}, \texttt{VERIFICATION\_RUN\_005}, \texttt{VERIFICATION\_RUN\_006}, and \texttt{VERIFICATION\_RUN\_007}. A permanent snapshot with a DOI is deposited on Zenodo.

\begin{acknowledgments}
This work was carried out within the Chyren Collaboration ($\odot\mathrm{\reflectbox{C}C}$), whose review process produced several of the retractions and rescopings recorded in the ledger.

\textbf{Funding:} No external funding. \textbf{Conflict of interest:} None.

\textbf{Author contributions and AI use:} R.W.Y.\ --- theoretical derivation, Lean formalization, numerical computation, and manuscript preparation. Large language model tools (Perplexity, Claude) contributed to draft generation, literature retrieval, and code review under the author's direction and audit. No AI system is an author, and the author takes full responsibility for all claims, derivations, and numerical results reported here. Every external citation in this manuscript was verified against its primary source.
\end{acknowledgments}

\bibliographystyle{apsrev4-2}
\bibliography{references}

\newpage
\appendix
\section*{Appendix A: Lean~4 Source Inventory and Verification Provenance}
\label{app:lean_theorems}

The tracked Lean~4 sources under \texttt{05\_lean\_formalization/} comprise \textbf{18 modules},
enumerated in Table~\ref{tab:lean_modules}. This appendix records the headline declarations per
module and the provenance of the verification run; it adds no claim beyond that table.

\paragraph{Verification provenance.}
O6 --- walked once in a clean worktree at 07185a6 (\texttt{lake exe cache get} + 17/17 PASS of the 17 targets then declared, \texttt{VERIFICATION\_RUN\_007}). Not yet demonstrated on a cold machine with empty host cache, and not yet a CI release gate. At that commit the gate \texttt{verify\_all\_proofs.sh} elaborated all 17 declared targets with exit status 0, no \texttt{sorry} in proof code, and axiom footprints confined to $\{\texttt{propext}, \texttt{Classical.choice}, \texttt{Quot.sound}\}$, under Lean \texttt{v4.33.0-rc1} and Mathlib \texttt{5eec30bc}. The transcript, per-target exit codes, toolchain and Mathlib revisions, and SHA-256 checksums are archived in \texttt{VERIFICATION\_RUN\_007/01\_lean/}.

\paragraph{What elaboration does and does not certify.}
A green gate run certifies that each statement follows from its encoded definitions. It does not
certify that those definitions model the physical system, and it does not certify assumptions
carried as typeclass fields, structure fields, or theorem hypotheses. The suite declares no Lean
global \texttt{axiom}; that is a statement about \texttt{axiom} declarations only. See
\texttt{THEORY\_ASSUMPTION\_AUDIT.md}.

\begin{table}[h!]
\centering
\small
\begin{tabularx}{\textwidth}{@{}l X@{}}
\toprule
\textbf{Module} & \textbf{Headline declarations} \\ \midrule
\multicolumn{2}{@{}l}{\emph{Proof modules}} \\ \midrule
\texttt{CosmologicalSector.lean} & \texttt{log\_two\_pos}, \texttt{log\_two\_lt\_one}, \texttt{matter\_density\_bounds}, \texttt{spatial\_flatness\_sum} \\
\texttt{CovariantCompletion.lean} & \texttt{raqual\_superluminal\_obstruction}, \texttt{disformal\_gamma\_ppn\_unity}, \texttt{preferred\_frame\_parameters\_zero}, \texttt{flrw\_projected\_gradient\_zero}, \texttt{no\_dynamical\_dark\_energy\_density} \\
\texttt{DeSitterExtremal.lean} & \texttt{desitter\_lapse\_horizon}, \texttt{expr\_cH\_over\_2pi\_pos}, \texttt{desitter\_flat\_limit}, \texttt{volume\_law\_weight\_nonneg} \\
\texttt{DualChannelDerivation.lean} & \texttt{dual\_channel\_flux\_algebra}, \texttt{mu\_derived\_inversion}, \texttt{mu\_derived\_deep\_mond\_upper\_bound}, \texttt{mu\_derived\_newtonian\_bound} \\
\texttt{GODActionKinematics.lean} & \texttt{dual\_channel\_poly\_identity}, \texttt{aqual\_simple\_mu\_ratio}, \texttt{btfr\_algebraic\_scaling} \\
\texttt{HorizonScale.lean} & \texttt{kms\_cancellation\_equilibrium}, \texttt{horizon\_acceleration\_ratio\_is\_one}, \texttt{verlinde\_entropic\_cancellation} \\
\texttt{ITActionClosure.lean} & \texttt{tauLaw\_eq\_simple\_mu\_poly}, \texttt{btfr\_deep\_mond}, \texttt{enclosed\_mass\_n2\_linear}, \texttt{flat\_rotation\_curve\_n2} \\
\texttt{MuProjection.lean} & \texttt{mu\_simple\_eq\_cos}, \texttt{quadratic\_law\_root\_unique}, \texttt{powerLaw\_solves\_dilaton\_eom}, \texttt{powerLaw\_iterated\_deriv}, \texttt{exp\_profile\_fails\_cubic} \\
\texttt{PPNLimits.lean} & \texttt{fractional\_deviation\_eq}, \texttt{solar\_system\_precision\_bound}, \texttt{cassini\_radar\_delay\_satisfied} \\
\texttt{RelativisticStability.lean} & \texttt{first\_derivative\_pos}, \texttt{second\_derivative\_pos}, \texttt{ghost\_free\_convexity} \\
\texttt{SOCasimirGenuine.lean} & \texttt{gen\_skew}, \texttt{sum\_dg}, \texttt{casimir\_defining\_rep}, \texttt{gen\_sq}, \texttt{casimir\_scalar\_eq} \\
\texttt{SkordisZlosnikEmbedding.lean} & \texttt{sz\_aqual\_reduction}, \texttt{dJ\_dY\_pos}, \texttt{sz\_tensor\_speed\_luminal}, \texttt{sz\_weak\_field\_lensing} \\
\texttt{SovereignRegularity.lean} & \texttt{chiral\_iff\_lipschitz\_constant}, \texttt{lipschitz\_implies\_angle\_modulus}, \texttt{bkm\_no\_blowup}, \texttt{sovereign\_regularity\_theorem} \\
\texttt{TensorSpeed.lean} & \texttt{c\_T\_sq\_at\_zero}, \texttt{maxwellian\_c13\_vanishes}, \texttt{foster\_jacobson\_alpha\_1\_eval}, \texttt{speed\_ratio\_unity\_iff}, \texttt{gw170817\_deviation\_of\_pos} \\
\texttt{YettParadigm.lean} & \texttt{ramanujan\_yett\_spectral\_gap\_pos}, \texttt{adccl\_trajectory\_bounded}, \texttt{chiral\_phase\_stable} \\ \midrule
\multicolumn{2}{@{}l}{\emph{Assumption declarations --- not physics results}} \\ \midrule
\texttt{AXIOMS\_V2.lean} & Typeclasses \texttt{Axiom\_A1}--\texttt{A5}; \texttt{derived\_simple\_mu\_bounds}, \texttt{deep\_mond\_baryonic\_scaling} \\ \midrule
\multicolumn{2}{@{}l}{\emph{Diagnostic helpers --- no proof content}} \\ \midrule
\texttt{PrintAxioms.lean} & \texttt{\#print axioms} for the \texttt{CovariantCompletion} results \\
\texttt{PrintAxiomsD8.lean} & \texttt{\#print axioms} for the tensor-speed results \\ \bottomrule
\end{tabularx}
\caption{Headline declarations in the 18 tracked Lean~4 source modules, enumerated from
\texttt{git ls-files}. Epistemic status per module is given in
Table~\ref{tab:lean_modules}; this table lists declaration names only.}
\end{table}

\paragraph{Scope note on \texttt{SovereignRegularity.lean}.}
The module bounds pointwise vorticity by an assumed control field. It does not formalize an
integral and does not formalize a partial differential equation, so it does not establish the
Beale--Kato--Majda criterion \cite{BealeKatoMajda1984}, which concerns
$\int_0^T \lVert \omega \rVert_\infty \, \mathrm{d}t$. \texttt{sovereign\_regularity\_theorem}
restates the module's own hypothesis.

\section*{Appendix B: SPARC Metric Definitions}
\label{app:sparc_metrics}

For each galaxy $g$ with $N_g$ observed kinematic radii $R_i$, observed velocities
$V_{\text{obs},i}$, and observational uncertainties $\sigma_{V,i}$, the model velocity is:
\begin{equation}
V_{\text{model}}(R_i) = V_{\text{bar}}(R_i) \cdot \left[ \frac{1}{2} + \sqrt{\frac{1}{4} + \frac{a_0}{g_{\text{bar}}(R_i)}} \right]^{1/2},
\end{equation}
where $g_{\text{bar}}(R_i) = V_{\text{bar}}(R_i)^2 / R_i$, with
$V_{\text{bar}}^2(R) = |V_{\text{gas}}|V_{\text{gas}} + \Upsilon_{\text{disk}} |V_{\text{disk}}|V_{\text{disk}} + \Upsilon_{\text{bulge}} |V_{\text{bulge}}|V_{\text{bulge}}$.

\subsection*{Per-galaxy metric}
\begin{equation}
\frac{\chi^2_{\text{data}}}{N_g} = \frac{1}{N_g} \sum_{i=1}^{N_g} \left( \frac{V_{\text{obs},i} - f_d V_{\text{model}}(f_d R_i)}{\sigma_{V,i}} \right)^2.
\end{equation}

These are definitions only. The sample size, parameter count, and all fitted values are carried
by the frozen artifacts named in Appendix~C; they are deliberately not restated here, so that
this appendix cannot drift from them.

\section*{Appendix C: Parameter Accounting and Cross-Validation}
\label{app:param_derivation}

\paragraph{Superseded accounting --- do not quote.}
Earlier revisions of this appendix reported an aggregate residual over
$N_{\text{data}} = 3{,}391$ points with $N_{\text{par}} = 382$ and
$\text{DOF}_{\text{nom}} = 3{,}009$, and a 5-fold cross-validation evaluated against
\emph{fixed} population priors. Both are superseded: items \textbf{D4.1} and \textbf{D4.2} of
\texttt{EPISTEMIC\_BOUNDARY\_v1.5.0.md} record the method defects --- radial points treated as
independent, and a cross-validation that held $a_0$ fixed across folds and so was not a
cross-validation. Those numbers are retained here only to be labelled, and must not be quoted
as current results.

\paragraph{Current accounting.}
The parameter ledger, sample definition, systematic budget, and honest per-fold
cross-validation are carried by the frozen artifacts on \texttt{main}:

\begin{itemize}
  \item \texttt{02\_galaxy\_dynamics/A0\_MEASUREMENT.json} --- working $a_0$ with the
        statistical and systematic budget (item D4.3).
  \item \texttt{02\_galaxy\_dynamics/PARAMETER\_LEDGER.json} --- tiered parameter counts and
        per-tier medians (items D4.7, D4.8).
  \item \texttt{02\_galaxy\_dynamics/NFW\_CONSTRAINED.json} --- concentration-prior comparison
        (item D4.9).
  \item \texttt{02\_galaxy\_dynamics/A0\_ESTIMATE.json} --- cross-validation retraining $a_0$
        per fold (item D4.6).
  \item \texttt{02\_galaxy\_dynamics/SPARC\_PARAMETER\_BUDGET.md} --- the budget derivation.
\end{itemize}

Where this appendix and \texttt{EPISTEMIC\_BOUNDARY\_v1.5.0.md} appear to differ, the boundary
document governs.

\section*{Appendix D: Reproduction Commands}
\label{app:reproduction_commands}

All paths are relative to the repository root. SPARC rotation-curve tables are not vendored;
see \texttt{02\_galaxy\_dynamics/SPARC\_DATA.md} for the official source and the checksum
manifest.

\begin{verbatim}
# 1. Lean 4 formal suite (all 17 declared targets)
cd 05_lean_formalization
lake exe cache get          # fresh-clone path (see O6)
./verify_all_proofs.sh      # exits 0 only if every target elaborates cleanly

# 2. SPARC measurement and parameter ledger
#    Requires the rotation-curve tables; see 02_galaxy_dynamics/SPARC_DATA.md
cd 02_galaxy_dynamics
python3 a0_measure.py
python3 parameter_ledger.py

# 3. Manuscript PDF
pdflatex -interaction=nonstopmode final_manuscript.tex
bibtex   final_manuscript
pdflatex -interaction=nonstopmode final_manuscript.tex
pdflatex -interaction=nonstopmode final_manuscript.tex
\end{verbatim}

\section*{Appendix E: Artifact Manifests}
\label{app:artifact_manifest}

Earlier revisions printed a table of truncated SHA-256 prefixes, which cannot be verified. The
full manifests are tracked in the repository and should be checked directly:

\begin{itemize}
  \item \texttt{VERIFICATION\_RUN\_001/02\_sparc\_strict\_135/RAW\_DATA\_MANIFEST.sha256} ---
        SHA-256 of each of the 175 \texttt{*\_rotmod.dat} inputs. Verify with
        \texttt{sha256sum -c} from the data directory.
  \item \texttt{VERIFICATION\_RUN\_003/01\_lean/SHA256SUMS.txt} --- checksums of the Lean gate
        transcript and its machine-readable result record.
  \item \texttt{05\_lean\_formalization/lake-manifest.json} --- pinned Mathlib revision and
        transitive dependency revisions.
  \item \texttt{05\_lean\_formalization/lean-toolchain} --- pinned Lean toolchain.
\end{itemize}


\end{document}